As large language models (LLMs) proliferate, future models will increasingly be trained on AI-generated data. Prior work shows that recursively training models on their own outputs leads to ``model collapse.'' This paper investigates the concept of a ``Minimum Viable Population'' (MVP) for an LLM ecosystem to prevent such collapse. We empirically simulate recursive training loops to determine whether increasing the population size of models ($N > 1$) mitigates ecosystem degradation. Our results show that simply increasing the population size of identical models sharing a data pool actually accelerates collapse compared to a single-model baseline. A homogeneous population of three models exhibits a more severe drop in unique sentence ratios and a faster convergence to a biased mode than a single model. However, introducing epistemic diversity---specifically through varying generation temperatures---slightly but consistently mitigates this effect. We demonstrate that an MVP for LLMs cannot be defined merely by the count of agents; it fundamentally requires active maintenance of epistemic diversity. Our findings highlight that models must possess sufficiently different priors or sampling strategies to counterbalance the collective bias and sustain a healthy information ecosystem.